% Native semantic reconstruction of source physical page 32 (printed page 31).
% Provenance: IEEE Std 1516-2025, Clause 3.2 conclusion, Clause 4, and Rule 1 opening.
\HLAStartSemanticPage{32}{31}
\fontsize{10.02pt}{11.566pt}\selectfont
\begingroup
\hbadness=10000\relax
\addfontfeatures{FakeStretch=1.035}
\vspace*{-5.79pt}

\begin{hlaacronymrows}
\HLAAcronymRow{acro-som}{SOM}{simulation object model}
\HLAAcronymRow{acro-tradt}{TRADT}{time representation abstract data type}
\HLAAcronymRow{acro-tso}{TSO}{timestamp order}
\HLAAcronymRow{acro-uri}{URI}{uniform resource identifier}
\HLAAcronymRow{acro-url}{URL}{uniform resource locator}
\HLAAcronymRow{acro-wsprc}{WSPRC}{Web Service Provider RTI Component}
\HLAAcronymRow{acro-xml}{XML}{extensible markup language}
\end{hlaacronymrows}
\vspace{15.40pt}

\HLAClauseHeading{clause-4}{4. Summary of the HLA rules}
\vspace{13.12pt}

The rules for federations are as follows:\par
\vspace{15.06pt}

\begin{hlasummaryrules}
\input{output/semantic/requirement-summary-rule-1}
\HLASummaryRule{summary-rule-2}{In a federation, all simulation-associated object instance representation shall be in the\linebreak federates, not in the RTI.}
\HLASummaryRule{summary-rule-3}{During a federation execution, all exchange of FOM data among joined federates shall occur\linebreak via the RTI.}
\HLASummaryRule{summary-rule-4}{During a federation execution, joined federates shall interact with the RTI in accordance with\linebreak the HLA interface specification.}
\HLASummaryRule{summary-rule-5}{During a federation execution, an instance attribute shall be owned by at most one joined\linebreak federate at any given time.}
\end{hlasummaryrules}
\vspace{9.03pt}

The rules for federates are as follows:\par
\vspace{15.06pt}

\begin{hlasummaryrules}
\setcounter{enumi}{5}
\input{output/semantic/requirement-summary-rule-6}
\HLASummaryRule{summary-rule-7}{Federates shall be able to update and/or reflect any instance attributes and send and/or receive\linebreak interactions, as specified in their SOMs.}
\HLASummaryRule{summary-rule-8}{Federates shall be able to transfer and/or accept ownership of instance attributes dynamically\linebreak during a federation execution, as specified in their SOMs.}
\HLASummaryRule{summary-rule-9}{Federates shall be able to vary the conditions (e.g., thresholds) under which they provide\linebreak updates of instance attributes, as specified in their SOMs.}
\HLASummaryRule{summary-rule-10}{Federates shall be able to manage local time in a way that allows them to coordinate data\linebreak exchange with other members of a federation.}
\end{hlasummaryrules}
\vspace{15.26pt}

\HLAClauseHeading{clause-5}{5. Federation rules}
\vspace{22.10pt}

\HLADetailedRuleHeading{detailed-rule-1}{5.1 Rule 1}
\vspace{10.78pt}

\HLANormativeRuleStatement{Federations shall have an HLA FOM, documented in accordance with the HLA OMT.}
\vspace{11.82pt}

All data to be exchanged in accordance with the HLA shall be documented in a FOM. The FOM is\linebreak
specified using one or more FOM modules and one MOM and Initialization Module (MIM). A FOM shall\linebreak
document the agreement among federates on data to be exchanged using the HLA services during\linebreak
federation execution and the minimal set of conditions of the data exchange (e.g., updates to be sent when\linebreak
changes exceed a certain value). As such, a FOM is an essential element in defining a federation. The HLA\linebreak
does not prescribe which data are included in a FOM (this task is the responsibility of the federation user\par
\endgroup
\clearpage
